Every government divides a limited budget among health, education, defence, social protection, transport and other functions. In practice the division is largely incremental: last year's allocation is adjusted for political priorities and expert judgment~\cite{davis1966,jones2005}. Machine learning (ML) can detect non-linear regularities in the socio-economic, fiscal and governance indicators that analysts cannot easily track jointly~\cite{athey2019,mullainathan2017}, and the allocations of comparable peer economies are rarely used systematically. Public-sector use, however, demands more than accuracy: a model output must be \emph{transparent}, and must remain an \emph{aid to, not a replacement for,} the accountable human decision-maker. This paper therefore asks a deliberately modest question: \emph{how much can machine learning add to last year's allocation, and what can peers and explanations tell us when there is no last year?}

\medskip\noindent\textbf{Use cases.} We see three defensible uses of such a model: (a) gap-filling and nowcasting of functional expenditure shares for statistical agencies and international organisations when official COFOG tables are late or missing; (b) peer benchmarking, i.e.\ showing what comparable economies typically spend; and (c) a benchmark dataset and protocol for research on compositional public-finance prediction. We do \emph{not} claim that the model tells a ministry what to allocate: the output is descriptive.

Existing ML work on public money mostly forecasts \emph{aggregate} expenditure as a time series~\cite{gur2025,dragulin2025}, predicts which projects or actors win competitive funds~\cite{caravaggio2025,distefano2025}, or relates spending to outcomes~\cite{valle2022,guerrero2022}. Studies of budgets by function or account group are few and mostly single-country~\cite{santschi2024,radulescu2015,nonato2026,mohith2026}, and explanation methods such as SHAP and LIME are rarely tested against each other or against known drivers. A further obstacle is data: no ready-made benchmark combines official functional spending with comparable macro-governance indicators across countries. Our own earlier version of this work (v1) used a semi-synthetic benchmark whose excellent scores turned out to be an artefact of the targets (Section~\ref{sec:audit}); the present paper replaces it with official data.

\medskip\noindent\textbf{Contributions.}
\begin{enumerate}
\item \textbf{An open, provenance-labelled dataset} (GovAlloc-COFOG~v1.0): \NMasterRows{} country-year rows for \NMasterCountries{} economies with \NFeatNum{} features, official COFOG functional allocations for \NCofogCountries{} economies (\NCofogRows{} complete country-years, 16 functions, Eurostat and OECD, cross-validated) and, with IMF data, ten COFOG divisions for \NTenEcon{} economies, plus official WDI sector series for \NWdiCountries{} economies; all legacy v1 columns are preserved.
\item \textbf{An evaluation protocol for compositional targets} with the baselines practitioners use (persistence, peer mean), temporal, leave-country-out and leave-region-out splits, a budget-reallocation metric (total-variation distance, TVD), cluster-bootstrap intervals and a Holm correction for the cold-start comparisons.
\item \textbf{A benchmark of learners and negative results:} boosted trees, a nested-NNLS stack, a nearest-economies predictor, Dirichlet regression, neural networks, econometric panel baselines, shock-aware and institutional features, evaluated against persistence and the income-group mean. Several components (a bounded-simplex projection and split-conformal intervals) are included as \emph{evaluated components}, not as proposed methods.
\item \textbf{A quantitative evaluation of SHAP and LIME} for this task (agreement, stability, fidelity, a deletion test and a synthetic ground-truth test).
\item \textbf{A forensic audit of our own earlier semi-synthetic benchmark} showing how $R^2>0.99$ arose from eleven sector columns that are exact multiples of one series.
\end{enumerate}
\medskip\noindent\textbf{Principal findings.} With official targets, last year's allocation is very hard to beat: a learned correction using same-year covariates reduces the share of budget to be moved from \ResTVDPersist{} to \ResTVDHyb{} percentage points for the pre-declared learner seed; the gain (\ResPairedGain{} pp) is small and not robust, averaging \AnnSeedMeanGain{} pp over five learner seeds with a 95\% interval that excludes zero for only \AnnSeedNExcl{} of \AnnSeedN{}. Without history, no learner significantly beats the income-group mean for the 16-function panel after multiple-comparison correction, whereas on the wider ten-division panel a stacked learner and a nearest-economies average each beat it by about 1.3 pp, and health and R\&D intensity are predictable in unseen economies ($R^2$ 0.63 and 0.60 vs.\ 0.31 and 0.25, although a plain ridge regression already reaches 0.50 and 0.53). SHAP and LIME agree on direction (top-5 sign agreement \AgrSignMean) more than on ranking (Spearman \AgrRhoMean), and attributions on correlated macro indicators are ambiguous. Errors are two to three times larger for lower-income economies. We argue that a clear statement of this competence map is more useful to analysts than a headline accuracy.

Throughout, the quantity predicted is the \emph{executed} (outturn) allocation of general-government expenditure across COFOG functions, the only quantity reported consistently across countries; it is not the approved budget appropriation, and we use ``allocation'' in this sense. The annual ``forecast'' is a \emph{nowcast}: it uses same-year covariates, which are realised values that are not known when a budget is drafted, whereas persistence uses only past information (Section~\ref{sec:method}).

The paper is organised as follows: Section~\ref{sec:related} reviews related work; Section~\ref{sec:data} describes the data; Section~\ref{sec:method} the models and protocol; Section~\ref{sec:results} the results; Section~\ref{sec:disc} discusses implications, threats to validity and governance; Section~\ref{sec:conc} concludes.
